% ===== 原文 pp.37–38 =====
% =====================================================================
%  tr05c —— 原文 pp.37–38（§4.3 问题集 III：问题 11–15）
%  编号覆盖：问题 11、12、13、14、15；小节标题 \subsection{问题集 III}(4.3)
%  起点：p37 开头（4.3 小节）；终点：p38 末（问题 15 末），p39 由 tr-06 从 §5 接
% =====================================================================

% [原文 p.37]
\subsection{问题集 III}

本问题集旨在使读者熟悉黎曼面上的作用与群上的作用，以及它们在 Higgs 丛模空间上诱导的作用。其中的问题兼收讲义中的问题与一些额外的问题。

\begin{prob}{11} \textbf{拓扑不变量。} $U(p,q)$-Higgs 丛的 $(\mathrm{B},\mathrm{A},\mathrm{A})$-膜中，与 $(W_1\oplus W_2,\Phi)$ 相联系的拓扑不变量是次数 $\deg(W_1)$ 与 $\deg(W_2)$。当 $p=q$ 时，由 [91]，这些不变量可以用 $S$ 上线丛 $L$ 的次数以及 $S$ 上那些使 $\sigma$ 在 $L$ 的纤维上的线性作用为 $-1$ 的分歧点数来表示。这些不变量满足若干性质，下面逐一来看：
% [Tutor Note: p.37 问题 11 的标题在原文 PDF 中印作 "textbfTopological invariants."，
%  系 LaTeX 源码漏写反斜杠（\textbf 未生效）所致；按体例恢复为粗体小标题“拓扑不变量。”，
%  此处照录原文状况备查。]
\begin{enumerate}[label=(\roman*),leftmargin=2.2em]
\item 利用 [9, Theorem 4.12] 中的 Lefschetz 不动点公式，说明 $L$ 的次数的奇偶性与 $\sigma$ 作用为 $-1$ 的点数必须相同。
\item 按照 [86]，与 $U(p,p)$-Higgs 丛相联系的 Toledo 不变量 $\tau(\deg(W_1),\deg(W_2))$ 定义为 $\tau(\deg(W_1),\deg(W_2)) := \deg(W_1) - \deg(W_2)$。利用习题 11 把该不变量表示为 $\sigma$ 的不动点的表达式，并得到自然的界。
\item ((*)) 对于具有极大 Toledo 不变量的 $\SU(p,p)$-Higgs 丛，谱数据由 $\Prym(S/\sigma,\Sigma)$ 的一个覆盖给出。对于 $\SU(p,p+1)$-Higgs 丛，可以用 cameral 数据得到谱数据，见 [83]。把 [91] 的工具与 [104] 的工具结合起来，给出 $\SU(p,1)$ 谱数据的描述。
% [Tutor Note: p.37 问题 11 (iii) 原文作 "maximal Toledo Toledo invariant"，"Toledo" 重复，
%  应为 "maximal Toledo invariant"；同条 "Use Exercise 11" 中的 "Exercise 11" 实指本问题集的
%  Problem 11（作者习惯混用 exercise/problem），译文照原文作“习题 11”。]
\end{enumerate}
\end{prob}

\begin{prob}{12} \textbf{来自有限群的 $(\mathrm{B},\mathrm{B},\mathrm{B})$-膜。} 回忆一下，亏格为 $g$ 的黎曼面 $\Sigma$ 上平坦 $\Sl(2,\C)$-联络的模空间维数为 $\dim\cM_g = 6g-6$；而对于亏格为 $\gamma$、带 $n$ 个孔、局部单值性固定（特征值为单根）的黎曼面，其上平坦 $\Sl(2,\C)$-联络的模空间 $\cM_{\gamma,n}$ 的维数为 $\dim\cM_{\gamma,n} = 6\gamma - 6 + 2n$，前提是 $\gamma \geqslant 2$，或 $\gamma \geqslant 1$ 且 $n \geqslant 1$，或 $\gamma = 0$ 且 $n \geqslant 4$。否则，$\cM_{\gamma,n}$ 的维数为 $0$ 或 $2$。
% [Tutor Note: p.37 问题 12 原文 "the dimension of the moduli space of flat space is SL(2,C)-connections"
%  一句中 "of flat space is" 为衍文（疑由 "of flat SL(2,C)-connections" 与 "space" 混排而成），
%  句意应为"平坦 $\Sl(2,C)$-联络的模空间维数为 $6g-6$"；译文按正意译出。]
\begin{enumerate}[label=(\roman*),leftmargin=2.2em]
\item 利用 Riemann--Hurwitz 公式，给出黎曼面 $\Sigma$ 与 $\Sigma/\Z_p$ 的亏格 $g$ 与 $\gamma$ 之间的关系。
\item 设 $g \geqslant 2$，$\Gamma$ 是阶为 $h$ 的有限群，以全纯自同构作用于 $\Sigma$。利用问题 12.i.，证明：若 $\Sigma$ 上等变平坦 $\Sl(2,\C)$-联络的模空间的某个分支具有 $\cM_g$ 维数的一半，则对某个自然数 $k \in \N$ 有 $h = 2^k$。
\item 对于问题 12.ii. 中 $g = 1$ 的情形，你能说些什么？
\end{enumerate}
\end{prob}

\begin{prob}{13} \textbf{等变与反等变谱数据。} 给定紧黎曼面 $\Sigma$，可以通过全纯截面 $s \in H^0(\Sigma,L^2)$ 构造 $\Sigma$ 的二重覆盖，其中 $L$ 是 $\Sigma$ 上的全纯线丛。若只考虑具有单零点的截面，则存在一个在 $s$ 的零点处分歧的二重覆盖 $\pi : S \to \Sigma$，使得存在平方根 $t \in H^0(S,\pi^*L)$ 满足 $t^2 = s$。

% [原文 p.38]
\begin{enumerate}[label=(\roman*),leftmargin=2.2em]
\item 在考虑 Hitchin 纤维化时，Hitchin 底中的光滑轨迹由具有单零点的 $H^0(\Sigma,K^2)$ 中截面给出。改取 $L = K^2$，证明存在（唯一的）二重覆盖 $\bar\pi : \bar S \to \Sigma$，它在 $s \in H^0(\Sigma,L^2)$ 的零点处分歧，且存在平方根 $t \in H^0(S,\pi^*L)$ 满足 $t^2 = s$。
% [Tutor Note: p.38 问题 13.i 原文的覆盖记作 $\bar\pi:\bar S\to\Sigma$，而下一句平方根所在的空间
%  写作 $H^0(S,\pi^*L)$（用 $S$ 与 $\pi$）；同一小问内记号 $\bar\pi/\pi$、$\bar S/S$ 混用，照原文排。]
\item 注意，对具有谱线丛 $L$ 的 $\Sl(2,\C)$-Higgs 场 $\Phi$ 而言，$\Phi$ 的谱线丛是 $\sigma^*L$，其中 $\sigma$ 是交换二重谱曲线两个叶片的自然对合。证明：关于 $\Sigma$ 上作用的无不动点对合 $\tau$ 的反等变 $\Sl(2,\C)$-Higgs 丛，其谱数据由
\[
    \cP^{\vee} := \{L \in \Jac(S) \mid \sigma^*L = L^*,\ \tau^*L = L^*\} \subset \Prym(S,\Sigma)
\]
确定。
\item 回顾：分裂实形式的 Higgs 丛对应于复 Higgs 丛在 Hitchin 纤维中的 2-挠点。据此对 $\Sl(2,\R)$-Higgs 丛证明与问题 13.ii. 类似的结果。
\end{enumerate}
\end{prob}

\begin{prob}{14} \textbf{来自实结构的 $(\mathrm{B},\mathrm{A},\mathrm{A})$-膜。} 我们已经看到实 Higgs 丛如何通过对李代数的对合来定义。这些作用也可用来理解该膜的谱数据描述以及标记膜的连通分支的拓扑不变量的性质。
\begin{enumerate}[label=(\roman*),leftmargin=2.2em]
\item 表 4.8 中相应的对合 $\theta$ 给出分解 $\fu = \fh \oplus \mathrm{i}\fm$。给出 $\fm$ 与 $\fh$ 以及实形式 $\fg = \fh\oplus\fm$ 的显式描述。
\item ((*)) 如 [83] 所示，$\Theta_{\mathrm{Sp}(2p,2p)}$ 的不动点集的非阿贝尔化也可以通过 cameral 覆盖来看出。用 cameral 覆盖刻画 $\sigma$ 的作用。
\item ((*)) 考虑 [1] 中“强实形式”（strong real form）的概念，描述相应的 Higgs 丛，并给出 $\Theta_G$ 的一个定义，使得不再出现上一段所述的问题。
\end{enumerate}
\end{prob}

\begin{prob}{15} \textbf{来自实结构的 $(\mathrm{A},\mathrm{B},\mathrm{A})$-膜。} 设 $\Sigma$ 亏格为 2，带有一个实结构 $f$，其不变量为 $(n,a) = (3,0)$。

考虑一个 $\Sl(2,\C)$-Higgs 丛，其谱曲线为 $S = \{0 = \eta^2 - q\}$，其中 $q$ 是具有单零点的二次微分。对 $z\in\Sigma$，对合 $f$ 诱导对合 $\tilde f(\eta,z) = (f^*(\eta), f(z))$。
\begin{enumerate}[label=(\roman*),leftmargin=2.2em]
\item 把 $q$ 换成 $\alpha q$（$\alpha\in\R_+$）会有什么效果？把 $q$ 换成 $-q$ 呢？
\item 设 $\Sigma$ 是超椭圆的，用 $\PP^1$ 中的两个点 $a_1,a_2$ 描述 $\Sigma$ 上的二次微分系统。
\item 问题 15.ii. 中的 $a_1$ 与 $a_2$ 须满足什么条件，才能得到满足 $f^*(\bar q) = q$ 的二次微分 $q$？
\end{enumerate}
\end{prob}
